\documentclass[12pt, reqno]{amsart}
\usepackage{latexsym, amssymb, amscd, xypic, amsthm, setspace}
\input xy
\xyoption{all}

%%: Font = times
\usepackage{mathptmx} % rm & math
\usepackage[scaled=0.90]{helvet} % ss
\usepackage{courier} % tt
\normalfont
\usepackage[T1]{fontenc}
\pagestyle{plain}


\newtheorem{theorem}{Theorem}[subsection]
\newtheorem{prop}[theorem]{Proposition}
\newtheorem{lemma}[theorem]{Lemma}
\newtheorem{corollary}[theorem]{Corollary}
\newtheorem{claim}[theorem]{Claim}

\theoremstyle{definition}
\newtheorem{de}[theorem]{Definition}
\newtheorem{notation}{Notation}
\newtheorem{example}[theorem]{Example}
\newtheorem{examples}{Examples}
\newtheorem{remark}[theorem]{Remark}
\newtheorem{remarks}{Remarks}
\newtheorem{question}{Question}


\newtheorem*{nthm}{Theorem}
\newtheorem*{nlem}{Lemma}
\newtheorem*{nprop}{Proposition}
\newtheorem*{ncor}{Corollary}
\newtheorem*{nclaim}{Claim}
\newtheorem*{nde}{Definition}
\newtheorem*{nex}{Example}

\newtheoremstyle{test}{3pt}{3pt}{\itshape}{\parindent}{\bfseries}{---}{  }{\thmnote{#3}} \theoremstyle{test}
\newtheorem*{mythm}{}


%: paragraphs

\newcommand{\point}{\vspace{3mm}\par \noindent (\refstepcounter{subsection}{\bf \thesubsection) } }

\newcommand{\tpoint}[1]{\vspace{3mm}\par \noindent \refstepcounter{subsubsection}{\bf \thesubsubsection.}
  {\em #1. ---} }
\newcommand{\epoint}[1]{\vspace{3mm}\par \noindent \refstepcounter{subsection}{\bf \thesubsection.}
  {\em #1.} }
\newcommand{\spoint}{\vspace{3mm}\par \noindent \refstepcounter{subsection}{\bf \thesubsection.} }

\newcommand{\bpoint}[1]{\vspace{3mm}\par \noindent \refstepcounter{subsection}{\bf \thesubsection.}
  {\bf #1.} }



%: numbering and margins

\numberwithin{equation}{subsection}

\setlength{\oddsidemargin}{0cm} \setlength{\evensidemargin}{0cm}
\setlength{\marginparwidth}{0in}
\setlength{\marginparsep}{0in}
\setlength{\marginparpush}{0in}
\setlength{\topmargin}{0in}
\setlength{\headheight}{0pt}
\setlength{\headsep}{0pt}
\setlength{\footskip}{.3in}
\setlength{\textheight}{9.2in}
\setlength{\textwidth}{6.5in}
\setlength{\parskip}{0.25pt}
\setlength{\parindent}{0.25in}

%: math shortcuts


\DeclareMathOperator{\inv}{inv}
\DeclareMathOperator{\out}{Out}
\DeclareMathOperator{\gal}{Gal}
\renewcommand{\mod}{\mathrm{\, mod \, }}
\renewcommand{\d}[1]{{}^{#1}}
\renewcommand{\b}[1]{\mathbf{{#1}}}
\newcommand{\C}{\mathbb{C}}
\newcommand{\be}[1]{\begin{eqnarray} \label{#1}}
\newcommand{\ee}{\end{eqnarray}}
\newcommand{\re}{\mathsf{Re}}
\newcommand{\im}{\mathsf{Im}}
\newcommand{\rr}{\rightarrow}

\newcommand{\n}{\mathbb{N}}
\newcommand{\zee}{\mathbb{Z}}
\newcommand{\R}{\mathbb{R}}
\newcommand{\pr}{\mathbb{P}}
\renewcommand{\gcd}[2]{\mathsf{gcd}(#1, #2)}
\newcommand{\ov}[1]{\overline{#1}}
\renewcommand{\Re}{\mathrm{Re}}
\renewcommand{\Im}{\mathrm{Im}}



\title{Problem Set 2 \\ MATH 411: Honors Complex Variables, Fall 2026}
\date{last compiled: \today}

\numberwithin{equation}{section}
\refstepcounter{section}

\begin{document} %%
\maketitle
\begin{center}
Suggested completion date: Oct 5, 2026
\end{center}



\begin{enumerate}


\item Give the terms of order $\leq 3$ in the power series (using the expressions for $\cos(T)$, $\sin(T)$, and $e^T$ given in lecture):
\begin{enumerate}
\item $\dfrac{\sin T}{\cos T}$;
\item $\dfrac{e^T}{\sin T}$.
\end{enumerate}

\item Given complex numbers $a_0,a_1,u_1,u_2$, define $a_n$ for $n\geq 2$ by
\[
a_n=u_1a_{n-1}+u_2a_{n-2}.
\]
If
\[
T^2-u_1T-u_2=(T-\alpha)(T-\alpha')
\]
with $\alpha\neq\alpha'$, show that
\[
a_n=A\alpha^n+B(\alpha')^n
\]
for suitable $A,B$. Find $A,B$ in terms of $a_0,a_1,\alpha,\alpha'$. Consider the power series
\[
F(T)=\sum_{n=0}^{\infty}a_nT^n.
\]
Show that it represents a rational function $f(T)/g(T)$, where $f(T)$ and $g(T)$ are polynomials with $g(T)\neq 0$. Give the partial fraction decomposition of this rational function.

\item Determine the radius of convergence of the following series:
\begin{enumerate}
\item $\displaystyle\sum n^2z^n$;
\item $\displaystyle\sum \frac{n!}{n}z^n$;
\item If $f=\sum a_nz^n$ has radius of convergence $r>0$, show that $\sum n^da_nz^n$ has the same radius of convergence for any positive integer $d$.
\end{enumerate}

\item Prove the following:
\begin{enumerate}
\item The power series $\sum nz^n$ has radius of convergence $1$ but does not converge at any point of the unit circle.
\item The power series $\sum \frac{z^n}{n^2}$ has radius of convergence $1$ and converges at every point of the unit circle.
\item The power series $\sum \frac{z^n}{n}$ converges at every point of the unit circle except $z=1$.
\end{enumerate}

\item Find the terms of order $\leq 3$ in the power series expansion of
\[
f(z)=\frac{z-2}{(z+3)(z+2)}
\]
at $z=1$.

\end{enumerate}


\end{document}
